% !TeX root = ../../Book3.tex
\section{Mellin 变换与 \texorpdfstring{\(\theta\)}{\AnalysisPdfTheta} 函数}
\label{b3:sec:2602}

Gamma 因子出现的原因不是人为修饰：对一个高斯项做 Mellin 积分，二次尺度自然把参数变为 \(s/2\)，并产生 \(\pi^{-s/2}\Gamma(s/2)\)。

\subsection{Mellin 变换}
\label{b3:subsec:260201}

记 \(\vartheta(t)=\sum_{n\in\mathbb Z}e^{-\pi n^2t}\)、\(h(t)=(\vartheta(t)-1)/2=\sum_{n\ge1}e^{-\pi n^2t}\)，\(t>0\)。对 \(\sigma>1\) 可用 Tonelli 估计绝对值并换序：
\[
\int_0^\infty h(t)t^{s/2-1}\,\mathrm{d}t
=\sum_{n\ge1}(\pi n^2)^{-s/2}\Gamma(s/2)
=\pi^{-s/2}\Gamma(s/2)\zeta(s).
\]
绝对积分为 \(\pi^{-\sigma/2}\Gamma(\sigma/2)\zeta(\sigma)<\infty\)，所以复参数换序并非仅由每一项可积推出。

\subsection{Gamma 函数}
\label{b3:subsec:260202}

Gamma 函数的无零点性使这个乘法不会在右半平面创造新零点。其极点以后却会抵消 \(\zeta\) 的负偶数零点，所以除去 Gamma 因子时必须单独查看局部阶数。特别是 \(s=0\)，Gamma 有极点而 \(\zeta\) 最终并不为零，不能把所有 Gamma 极点一概称为平凡零点的位置。

\subsection{\texorpdfstring{\(\theta\)}{\AnalysisPdfTheta} 函数}
\label{b3:subsec:260203}

由\secref{b3:sec:2401}的高斯求和，\(\vartheta(t)=t^{-1/2}\vartheta(1/t)\)。于是
\[
h(t)=t^{-1/2}h(1/t)+\frac12(t^{-1/2}-1).
\]
这是延拓的关键等式：右侧第一项在 \(t\downarrow0\) 时带指数小因子，另外两项完整记录原点的幂次奇性。只减去常数一还不能使原点积分对所有参数收敛；还必须处理 \(t^{-1/2}\) 项。

\subsection{热核观点}
\label{b3:subsec:260204}

函数 \(\vartheta(t)\) 也可视为圆周高斯热核在原点的值，\(n^2\) 对应离散频率的平方。减去一是在去掉零频率。小时间的 \(t^{-1/2}\) 主项反映一维局部尺度；Mellin 积分把这个幂次转成 \(s=1\) 的极点。这个解释无需将谱语言当作证明，全部常数已由高斯换元确定。

\ProseParagraphBreak

为了使热核的归一化完全明确，在周期一的圆周上取 Fourier 模态 \(e^{2\pi\mathrm inx}\)。算子 \(-\partial_x^2\) 的对应特征值为 \(4\pi^2n^2\)，故热核在时间 \(\tau\) 的对角值是
\[
 K_\tau(0)=\sum_{n\in\mathbb Z}e^{-4\pi^2n^2\tau}
          =\vartheta(4\pi\tau).
\]
令 \(t=4\pi\tau\)，才得到前面使用的 \(\vartheta(t)\)。时间参数若不同时缩放，会使 Mellin 变换中的 \(\pi\) 因子发生变化。对角热核的大小与变换公式都可以从 Fourier 级数逐项验证，无需额外假定抽象谱定理。

\ProseParagraphBreak

小时间的余项也有直接界。对 \(u\ge1\)，利用 \(n^2\ge1+3(n-1)\) 可得
\[
 0<h(u)\le\frac{e^{-\pi u}}{1-e^{-3\pi u}}
          \le\frac{e^{-\pi u}}{1-e^{-3\pi}}.
\]
代入参数反演，便有
\[
 \vartheta(t)=t^{-1/2}
       +O\bigl(t^{-1/2}e^{-\pi/t}\bigr),\qquad0<t\le1.
\]
这比仅说余项趋于零强得多。对任意固定的复参数紧集，指数小量乘上 \(t\) 的任意固定负幂仍可积；因此扣去主项后能够穿过任意有限条竖直线。下一节的全平面延拓正是把这一估计写成一条统一的积分公式。
